\documentclass[11pt,a4paper]{article}
\usepackage[T1]{fontenc}
\usepackage[utf8]{inputenc}
\usepackage[french]{babel}
\usepackage{lmodern}
\usepackage{textcomp}
\usepackage{amsmath,amssymb}
\usepackage{geometry}
\usepackage{tabularx,array}
\usepackage[table]{xcolor}
\usepackage{tikz}
\geometry{left=8mm,right=8mm,top=7mm,bottom=7mm}
\setlength{\parindent}{0pt}
\pagestyle{empty}
\setlength{\tabcolsep}{3pt}
\renewcommand{\arraystretch}{1.05}
\renewcommand{\tabularxcolumn}[1]{m{#1}}

% --- Palette pensée pour une impression noir et blanc ---
\definecolor{canlight}{gray}{0.93}
\definecolor{canmid}{gray}{0.72}
\definecolor{candark}{gray}{0.22}
\definecolor{answerblue}{RGB}{0,82,165}

% --- Colonnes ---
\newcolumntype{Q}{>{\centering\arraybackslash}X}
\newcolumntype{N}{>{\centering\arraybackslash}m{6mm}}
\newcolumntype{R}{>{\centering\arraybackslash}m{22mm}}

% --- En-tête ---
\newcommand{\canheader}[1]{%
  \begin{center}
    {\LARGE\bfseries COURSE AUX NOMBRES}\par
    \vspace{0.7mm}
    {\color{canmid}\rule{27mm}{0.45pt}\hspace{2mm}\color{black}$\diamond$\hspace{1.4mm}{\large\bfseries #1}\hspace{1.4mm}$\diamond$\hspace{2mm}\color{canmid}\rule{27mm}{0.45pt}\color{black}}
  \end{center}
  \vspace{0.8mm}
  \renewcommand{\arraystretch}{1.25}%
  \begin{tabularx}{\linewidth}{>{\centering\arraybackslash}X >{\centering\arraybackslash}X >{\centering\arraybackslash}m{40mm}}
    \rowcolor{canlight}
    \footnotesize\textbf{30 questions} -- 8 minutes &
    \footnotesize\textbf{Sans calculatrice} &
    \footnotesize\textbf{Score :} \rule{13mm}{0.45pt}\,/30
  \end{tabularx}
  \renewcommand{\arraystretch}{1.05}%
  \vspace{1mm}
  \begin{center}
    {\scriptsize Réponse uniquement \quad$\diamond$\quad Reporter le score sur la fiche annuelle.}
  \end{center}
  \vspace{1.2mm}
}

\newcommand{\qrow}[2]{\textbf{#1} & \parbox[c][15.0mm][c]{\linewidth}{\centering #2} & \\ \hline}
\newcommand{\qrowfig}[2]{\textbf{#1} & \parbox[c][27mm][c]{\linewidth}{\centering #2} & \\ \hline}
\newcommand{\qrowcorr}[3]{\textbf{#1} & \parbox[c][15.0mm][c]{\linewidth}{\centering #2} & {\color{answerblue}\bfseries #3} \\ \hline}
\newcommand{\qrowfigcorr}[3]{\textbf{#1} & \parbox[c][27mm][c]{\linewidth}{\centering #2} & {\color{answerblue}\bfseries #3} \\ \hline}

\newcommand{\cantable}[1]{%
\arrayrulecolor{canmid}
\begin{tabularx}{\linewidth}{N|Q|R}
\rowcolor{canlight}\textbf{N\textsuperscript{o}}&\centering\textbf{Question}&\textbf{Réponse}\\\hline
#1
\end{tabularx}
\arrayrulecolor{black}
}

\newcommand{\hypotfigure}{%
\begin{tikzpicture}[baseline=(current bounding box.center),scale=0.72,line width=0.6pt]
  \node[font=\bfseries\normalsize] at (1.6,2.55) {Quelle est l'hypoténuse ?};
  \coordinate (A) at (0,0); \coordinate (B) at (0,1.85); \coordinate (C) at (3.2,0);
  \draw (A)--(B)--(C)--cycle;
  \draw (0,0) rectangle (0.32,0.32);
  \node[below left] at (A) {\small A};
  \node[above left] at (B) {\small B};
  \node[below right] at (C) {\small C};
\end{tikzpicture}}


\newcommand{\canheadercorr}[1]{%
  \begin{center}
    {\LARGE\bfseries COURSE AUX NOMBRES}\par
    \vspace{0.7mm}
    {\color{canmid}\rule{27mm}{0.45pt}\hspace{2mm}\color{black}$\diamond$\hspace{1.4mm}{\large\bfseries #1}\hspace{1.4mm}$\diamond$\hspace{2mm}\color{canmid}\rule{27mm}{0.45pt}\color{black}}
  \end{center}
  \vspace{0.8mm}
  \renewcommand{\arraystretch}{1.25}%
  \begin{tabularx}{\linewidth}{>{\centering\arraybackslash}X >{\centering\arraybackslash}X >{\centering\arraybackslash}m{40mm}}
    \rowcolor{canlight}
    \footnotesize\textbf{30 questions} -- 8 minutes &
    \footnotesize\textbf{Sans calculatrice} &
    \footnotesize\textcolor{answerblue}{\textbf{CORRIGÉ}}
  \end{tabularx}
  \renewcommand{\arraystretch}{1.05}%
  \vspace{1mm}
  \begin{center}
    {\scriptsize Les réponses apparaissent en \textcolor{answerblue}{\textbf{bleu}}.}
  \end{center}
  \vspace{1.2mm}
}


\begin{document}
\canheader{CAN 3}

\begin{minipage}[t]{0.493\textwidth}
\normalsize
\cantable{
\qrow{1}{$9\times6$}
\qrow{2}{$100-37$}
\qrow{3}{$-3+8-2$}
\qrow{4}{Compléter : $56\div\ldots=8$}
\qrow{5}{$50\,\%$ de 74}
\qrow{6}{La somme du produit de 4 par 5 et de la différence entre 12 et 7}
\qrow{7}{$0{,}4+0{,}35$}
\qrow{8}{$\dfrac57$ multiplié par quel nombre donne 5 ?}
\qrow{9}{$0{,}006\times1000$}
\qrow{10}{Dans une série de 28 valeurs, une fréquence de $0{,}25$ correspond à quel effectif ?}
\qrow{11}{Le plus petit : $3{,}09$ ; $3{,}9$ ; $3{,}099$}
\qrow{12}{L'opposé de $8$}
\qrow{13}{$\dfrac34$ correspond à quel pourcentage ?}
\qrow{14}{2 h 45 min = $\ldots$ min}
\qrow{15}{$20-3\times4$}
}
\end{minipage}\hfill
\begin{minipage}[t]{0.493\textwidth}
\normalsize
\cantable{
\qrow{16}{$\dfrac12+\dfrac14$}
\qrow{17}{Périmètre d'un carré de côté $6{,}5$ cm}
\qrow{18}{Moyenne de 5 ; 10 ; 15 ; 10}
\qrow{19}{Quel nombre vaut $3\times100+4\times10+7\times0{,}1$ ?}
\qrow{20}{Un volume de $27\,\mathrm{cm}^3$ correspond à combien de cubes de 1 cm de côté ?}
\qrow{21}{315 est-il divisible par 9 ?}
\qrow{22}{$17\times98+17\times2$}
\qrow{23}{$2{,}4$ kg = $\ldots$ g}
\qrow{24}{6 croissants coûtent 9 \texteuro. Quel est le prix d'un croissant ?}
\qrow{25}{Dans un triangle, deux angles mesurent $90^\circ$ et $28^\circ$. Le troisième mesure :}
\qrow{26}{Un carré a un périmètre de 36 cm. Son côté mesure :}
\qrow{27}{Quel nombre, multiplié par lui-même, donne 81 ?}
\qrow{28}{Encadrer $-2{,}4$ par deux entiers consécutifs.}
\qrow{29}{Pour $x=2$, calculer $3x+1$.}
\qrow{30}{Réduire : $4x+2x$}
}
\end{minipage}

\clearpage

\canheadercorr{CAN 3}

\begin{minipage}[t]{0.493\textwidth}
\normalsize
\cantable{
\qrowcorr{1}{$9\times6$}{$54$}
\qrowcorr{2}{$100-37$}{$63$}
\qrowcorr{3}{$-3+8-2$}{$3$}
\qrowcorr{4}{Compléter : $56\div\ldots=8$}{$7$}
\qrowcorr{5}{$50\,\%$ de 74}{$37$}
\qrowcorr{6}{La somme du produit de 4 par 5 et de la différence entre 12 et 7}{$25$}
\qrowcorr{7}{$0{,}4+0{,}35$}{$0{,}75$}
\qrowcorr{8}{$\dfrac57$ multiplié par quel nombre donne 5 ?}{$7$}
\qrowcorr{9}{$0{,}006\times1000$}{$6$}
\qrowcorr{10}{Dans une série de 28 valeurs, une fréquence de $0{,}25$ correspond à quel effectif ?}{$7$}
\qrowcorr{11}{Le plus petit : $3{,}09$ ; $3{,}9$ ; $3{,}099$}{$3{,}09$}
\qrowcorr{12}{L'opposé de $8$}{$-8$}
\qrowcorr{13}{$\dfrac34$ correspond à quel pourcentage ?}{$75\,\%$}
\qrowcorr{14}{2 h 45 min = $\ldots$ min}{165 min}
\qrowcorr{15}{$20-3\times4$}{$8$}
}
\end{minipage}\hfill
\begin{minipage}[t]{0.493\textwidth}
\normalsize
\cantable{
\qrowcorr{16}{$\dfrac12+\dfrac14$}{$\dfrac34$}
\qrowcorr{17}{Périmètre d'un carré de côté $6{,}5$ cm}{26 cm}
\qrowcorr{18}{Moyenne de 5 ; 10 ; 15 ; 10}{$10$}
\qrowcorr{19}{Quel nombre vaut $3\times100+4\times10+7\times0{,}1$ ?}{$340{,}7$}
\qrowcorr{20}{Un volume de $27\,\mathrm{cm}^3$ correspond à combien de cubes de 1 cm de côté ?}{$27$}
\qrowcorr{21}{315 est-il divisible par 9 ?}{oui}
\qrowcorr{22}{$17\times98+17\times2$}{$1700$}
\qrowcorr{23}{$2{,}4$ kg = $\ldots$ g}{$2400$ g}
\qrowcorr{24}{6 croissants coûtent 9 \texteuro. Quel est le prix d'un croissant ?}{1,50 \texteuro}
\qrowcorr{25}{Dans un triangle, deux angles mesurent $90^\circ$ et $28^\circ$. Le troisième mesure :}{$62^\circ$}
\qrowcorr{26}{Un carré a un périmètre de 36 cm. Son côté mesure :}{9 cm}
\qrowcorr{27}{Quel nombre, multiplié par lui-même, donne 81 ?}{$9$}
\qrowcorr{28}{Encadrer $-2{,}4$ par deux entiers consécutifs.}{$-3<-2{,}4<-2$}
\qrowcorr{29}{Pour $x=2$, calculer $3x+1$.}{$7$}
\qrowcorr{30}{Réduire : $4x+2x$}{$6x$}
}
\end{minipage}
\end{document}
